% =====================================================================
% Module 1: Thu thap du lieu (Crawl Layer)
% ---------------------------------------------------------------------
% Cach su dung:
%   1. Dat file .tex nay cung thu muc voi bao cao chinh, hoac \input{} no.
%   2. Mo file "module1_diagrams.html" bang trinh duyet, chup man hinh tung
%      bieu do (Hinh 1.1 -> 1.5) va luoi thanh anh (PNG/PDF) dat cung thu muc
%      voi ten: module1_hinh1_1.png ... module1_hinh1_5.png
%   3. Go bo khoi "lstlisting setup" ben duoi neu bao cao da co package
%      "listings" va dinh ngha kieu Python.
% =====================================================================

% ---- Vietnamese encoding (xoa neu da co trong preamble) -------------
\usepackage[utf8]{inputenc}
\usepackage[T5]{fontenc}
% ---- lstlisting setup (xoa neu da co trong preamble) ----------------
\usepackage{listings}
\usepackage{xcolor}
\lstdefinestyle{pystyle}{
    language=Python,
    basicstyle=\ttfamily\footnotesize,
    keywordstyle=\color{blue}\bfseries,
    stringstyle=\color{teal},
    commentstyle=\color{gray}\itshape,
    showstringspaces=false,
    breaklines=true,
    frame=single,
    rulecolor=\color{lightgray},
    backgroundcolor=\color{white},
    tabsize=4,
    columns=fullflexible,
    keepspaces=true,
}
\lstset{style=pystyle}
% ---------------------------------------------------------------------

\usepackage{graphicx}
\usepackage{booktabs}
\usepackage{longtable}
\usepackage{hyperref}


\section{Module 1: Thu thập dữ liệu (Crawl Layer)}
\label{sec:crawl-layer}

Crawl Layer là stage đầu tiên của pipeline Lakehouse-Lite, chịu trách nhiệm thu
thập dữ liệu thô từ các website tuyển dụng và ghi ra file JSONL tạm cục bộ.
Nguyên tắc cốt lõi: \textbf{crawl layer chỉ ghi ra temp file, không bao giờ
talking tới S3/Supabase trực tiếp}. Việc upload lên Bronze là trách nhiệm của
storage layer ở các stage phía sau.


\subsection{Tổng quan kiến trúc}
\label{subsec:crawl-overview}

Crawl layer nhận trigger từ CLI hoặc Airflow DAG, chạy luồng orchestration
\texttt{crawl()} để phân trang, dedup URL, cào chi tiết và lưu từng trang ra
temp file ngay lập tức (streaming per-page save). Sơ đồ luồng tổng thể được
trình bày ở Hình~\ref{fig:crawl-pipeline}.

\begin{figure}[h!]
    \centering
    \includegraphics[width=\textwidth]{module1_hinh1_1.png}
    \caption{Luồng tổng thể của Crawl Layer: từ trigger đến lưu trữ Gold.
    Crawl layer chỉ ghi temp JSONL, các stage sau đọc temp.}
    \label{fig:crawl-pipeline}
\end{figure}

Pipeline chi tiết chạy theo thứ tự:
\begin{enumerate}
    \item \textbf{Crawl} -- ghi nối tiếp (append-only) ra
          \texttt{src/crawl\_layer/temp\_data/<source>\_jobs\_YYYYMMDD.jsonl}.
          Chạy lại cùng ngày sẽ cộng dồn dữ liệu.
    \item \textbf{Validate (tùy chọn)} -- kiểm tra dữ liệu temp trước khi upload.
    \item \textbf{Bronze} -- gzip temp JSONL, upload lên S3, rồi \textbf{xóa}
          temp của source đó (destructive).
    \item \textbf{Silver} -- làm sạch và xuất parquet theo partition ngày.
    \item \textbf{Supabase / MotherDuck (Gold)} -- tải silver lên serving layer.
\end{enumerate}


\subsection{Cấu trúc module crawler}
\label{subsec:crawl-structure}

Mỗi site là một package con trong \texttt{src/crawl\_layer/crawler/}. Cấu trúc
chuẩn được trình bày ở Hình~\ref{fig:crawl-package}.

\begin{figure}[h!]
    \centering
    \includegraphics[width=0.82\textwidth]{module1_hinh1_2.png}
    \caption{Cấu trúc package chuẩn của một module crawler.
    \texttt{http\_client.py} và \texttt{browser.py} là hai lựa chọn thay thế
    tùy theo stack kỹ thuật.}
    \label{fig:crawl-package}
\end{figure}

Một package crawler gồm các thành phần sau:

\begin{itemize}
    \item \texttt{\_\_init\_\_.py} -- rỗng, chỉ đánh dấu package.
    \item \texttt{\_\_main\_\_.py} -- điểm vào CLI: \texttt{argparse} +
          \texttt{asyncio.run} (kèm patch Windows cho nodriver).
    \item \texttt{config.py} -- hằng số tĩnh: URL, selector, timeout, tên source.
    \item \texttt{crawler.py} -- class \texttt{<Site>Crawler}: orchestration flow.
    \item \texttt{parser.py} -- class \texttt{<Site>Parser}: chuyển HTML thành
          \texttt{JobItem}.
    \item \texttt{http\_client.py} (stack aiohttp) hoặc \texttt{browser.py}
          (stack nodriver) -- tầng vận chuyển.
    \item \texttt{utils.py} -- helper nhỏ (encode keyword, sanitize text).
\end{itemize}

\textbf{Lưu ý:} không bao giờ có \texttt{\_\_init\_\_.py} ở cấp \texttt{src/} --
toàn bộ import dùng absolute \texttt{src.*}, phải \texttt{cd} về repo root trước
khi chạy \texttt{python -m}.


\subsection{Mô hình dữ liệu (Data Model)}
\label{subsec:crawl-datamodel}

Toàn bộ crawler chia sẻ dataclass gốc \texttt{JobItem} chứa các trường chung,
mỗi site kế thừa và chỉ thêm các trường đặc thù.

\begin{lstlisting}[caption={Dataclass gốc JobItem và các subclass per-site (src/crawl\_layer/data\_model/data\_class.py)},label={lst:jobitem}]
@dataclass
class JobItem:
    job_title: Optional[str] = None
    company_name: Optional[str] = None
    location: str | None = None
    job_industry: str | None = None
    job_description: str | None = None
    source_site: str | None = None
    job_url: str | None = None
    search_keyword: str | None = None
    scraped_at: str | None = None
    salary: Optional[str] = None
    benefits: str | None = None
    requirements: str | None = None

@dataclass
class TopCVJobItem(JobItem):
    company_size: Optional[str] = None
    job_type: str | None = None
    experience_level: str | None = None
    education_level: str | None = None
    job_position: str | None = None
    job_deadline: str | None = None

@dataclass
class ITViecJobItem(JobItem):
    company_size: Optional[str] = None
\end{lstlisting}

Quy tắc: chỉ thêm trường mà site \textbf{thực sự có}; mọi trường mặc định
\texttt{None} để parser có thể bỏ qua trường thiếu; tên trường phải khớp với
schema Silver ở hạ lưu.


\subsection{Cấu hình tĩnh (config.py)}
\label{subsec:crawl-config}

Mọi hằng số (URL, CSS/XPath selector, timeout, cookie, profile trình duyệt)
phải tập trung trong \texttt{config.py} để parser, http/browser và orchestrator
không bao giờ lệch nhau.

\begin{lstlisting}[caption={Ví dụ config.py cho stack nodriver (src/crawl\_layer/crawler/itviec/config.py)},label={lst:itviec-config}]
from src.storage_layer.MinIO_S3.config.path import DEFAULT_ENTITY_NAME

SOURCE_NAME = "itviec"
ENTITY_NAME = DEFAULT_ENTITY_NAME            # = "jobs"
BASE_URL = "https://itviec.com"
LOGIN_URL = "https://itviec.com/sign_in"
USERNAME_ENV = "ITVIEC_USERNAME"             # ten env var trong .env
PASSWORD_ENV = "ITVIEC_PASSWORD"

# Selectors - copy y trang that, class name co the bi hash chong bot
JOB_CARD_SELECTOR = "h3[data-url*='/it-jobs/']"
PREVIEW_PANEL_SELECTOR = "div[class*='preview-job-wrapper']"
NEXT_PAGE_SELECTOR = "a[rel='next']"

# Timing
LOGIN_TIMEOUT = 20.0
PAGINATION_DELAY_RANGE = (2.0, 4.0)

# Browser stealth args
BROWSER_ARGS = ("--no-sandbox",
                "--disable-blink-features=AutomationControlled", ...)
\end{lstlisting}

\texttt{SOURCE\_NAME} phải khớp chính xác với tiền tố file temp
(\texttt{<source>\_jobs\_*.jsonl}) và với \texttt{bronze\_source} trong cấu hình
Airflow. Không hardcode secret -- chỉ trỏ tới tên env var.


\subsection{Luồng orchestration (crawler.py)}
\label{subsec:crawl-flow}

\texttt{crawler.py} chỉ chứa luồng high-level, không chứa chi tiết fetch/parse.
Sơ đồ luồng \texttt{crawl()} được trình bày ở Hình~\ref{fig:crawl-flow}.

\begin{figure}[h!]
    \centering
    \includegraphics[width=0.7\textwidth]{module1_hinh1_4.png}
    \caption{Luồng orchestration của \texttt{Crawler.crawl()}: streaming
    per-page save để không mất dữ liệu khi crash giữa chừng.}
    \label{fig:crawl-flow}
\end{figure}

\begin{lstlisting}[caption={Pattern chung của class Crawler (rút gọn từ TopcvCrawler / ItviecCrawler)},label={lst:crawler-pattern}]
class <Site>Crawler:
    def __init__(self, keyword, max_pages, ...):
        self.keyword = keyword
        self.max_pages = max_pages
        self.<transport> = <Site>HttpClient(...) or <Site>Browser(...)
        self.parser = <Site>Parser()
        self._seen_urls: set[str] = set()      # dedup URL giua cac trang

    async def crawl(self) -> list[<Site>JobItem]:
        items: list[<Site>JobItem] = []
        async with self.<transport>:           # mo/dong session hoac browser
            for page_num in range(1, self.max_pages + 1):
                temp_items: list[<Site>JobItem] = []
                page_items = await self._scrape_current_page()
                items.extend(page_items)
                temp_items.extend(page_items)
                self._flush_batch(temp_items, page_num)  # ghi temp NGAY
                # ... next page ...
        return items

    def _flush_batch(self, page_items, page_num) -> None:
        if not page_items:
            return
        save_to_temp([asdict(i) for i in page_items], SOURCE_NAME, ENTITY_NAME)
\end{lstlisting}

Hai điểm bắt buộc trong luồng này:
\begin{enumerate}
    \item \textbf{Dedup URL} bằng \texttt{self.\_seen\_urls} để không cào trùng
          khi các trang overlap.
    \item \textbf{Streaming per-page save} qua \texttt{\_flush\_batch()}: ghi ra
          temp sau mỗi trang, nên nếu crash ở trang 5 thì trang 1--4 vẫn an toàn
          trên đĩa và bộ nhớ luôn bị chặn (bounded).
\end{enumerate}


\subsection{Hai stack kỹ thuật: aiohttp và nodriver}
\label{subsec:crawl-stacks}

Dự án dùng hai stack khác nhau tùy độ ``khó'' của site. So sánh chi tiết được
trình bày ở Hình~\ref{fig:crawl-stacks} và Bảng~\ref{tab:stack-compare}.

\begin{figure}[h!]
    \centering
    \includegraphics[width=\textwidth]{module1_hinh1_3.png}
    \caption{So sánh hai stack kỹ thuật: aiohttp (nhanh, không browser) và
    nodriver (chống bot, cần login).}
    \label{fig:crawl-stacks}
\end{figure}

\begin{table}[h!]
\centering
\caption{So sánh hai stack kỹ thuật thu thập dữ liệu}
\label{tab:stack-compare}
\begin{tabular}{p{0.28\textwidth} p{0.33\textwidth} p{0.33\textwidth}}
\toprule
\textbf{Tiêu chí} & \textbf{aiohttp (TopCV)} & \textbf{nodriver (ITviec / VietnamWorks)} \\
\midrule
File transport    & \texttt{http\_client.py} & \texttt{browser.py} \\
Cần login?        & Không & Có (Cloudflare + cookie) \\
Concurrency       & Cao -- \texttt{asyncio.gather} + semaphore & Thấp -- 1 session tuần tự \\
Thư viện          & \texttt{curl\_cffi} giả lập JA3 & \texttt{nodriver} headless Chrome \\
Số pha cào        & 2 pha (URL search → detail song song) & 1 pha (click card → snapshot panel) \\
Chống 403 CF      & Cơ chế Prober (Hình~\ref{fig:prober}) & Clearance gắn UA+IP+cookie \\
Linux/Docker cần  & Không & \texttt{xvfb-run} \\
Windows patch     & Không & Có (ProactorEventLoop) \\
\bottomrule
\end{tabular}
\end{table}

\subsubsection{Stack aiohttp -- HttpClient}
\label{subsec:aiohttp}

\texttt{TopcvHttpClient} sở hữu session \texttt{curl\_cffi}, semaphore giới hạn
concurrency và logic retry/backoff. Phương thức cốt lõi \texttt{fetch()} xử lý:
\texttt{BLOCK\_STATUS} (403) bằng cơ chế ``prober'' -- request đầu tiên phát hiện
block sẽ giăng \texttt{asyncio.Event} chặn toàn bộ request khác, ngủ 62 giây để
Cloudflare clearance tái lập, rồi mở rào khi thành công;
\texttt{RETRY\_STATUS} (429/5xx) bằng exponential backoff. Cơ chế prober được
trực quan ở Hình~\ref{fig:prober}.

\begin{figure}[h!]
    \centering
    \includegraphics[width=0.85\textwidth]{module1_hinh1_5.png}
    \caption{Cơ chế Prober trong \texttt{TopcvHttpClient.fetch()}: chỉ 1 request
    ``thám báo'' chạm Cloudflare, các request khác đứng chờ để tránh 429/ban IP.}
    \label{fig:prober}
\end{figure}

\begin{lstlisting}[caption={Logic fetch() rút gọn (src/crawl\_layer/crawler/topcv/http\_client.py)},label={lst:fetch}]
async def fetch(self, url, referer=None) -> str | None:
    async with self._semaphore:
        await asyncio.sleep(random.uniform(*self.request_delay))
        for attempt in range(1, self.max_retries + 1):
            if not is_prober:
                await self._is_clear.wait()        # cho neu dang GLOBAL PAUSE
            resp = await self.session.get(url, headers=headers)
            status = resp.status_code
            if status in BLOCK_STATUS:             # 403 -> Cloudflare
                # request dau tien lam prober, giang Event chan cac req khac
                ...
                await asyncio.sleep(62.0)          # cho CF clearance
            elif status in RETRY_STATUS:           # 429/5xx
                sleep_for = min(120.0, 2.0 * (2 ** (attempt - 1)))
                await asyncio.sleep(sleep_for + random.uniform(0, 1))
            elif 200 <= status < 300:
                if is_prober: self._is_clear.set() # mo rao cho cac req khac
                return resp.text
\end{lstlisting}

\subsubsection{Stack nodriver -- Browser}
\label{subsec:nodriver}

\texttt{ItviecBrowser} quản lý headless Chrome qua \texttt{nodriver}, có các
method: \texttt{login()} (điền form, chờ marker, raise \texttt{ItviecLoginError}
khi fail), \texttt{open\_search()}, \texttt{iter\_job\_panels()} (async generator
yielding \texttt{(job\_url, panel\_html)}), và \texttt{go\_to\_next\_page()}.
Lý do ITviec dùng 1 session tuần tự: Cloudflare gắn clearance với
(UA + IP + cookie đăng nhập), nên concurrency sẽ tự đánh bại mục đích.


\subsection{Parser và utils}
\label{subsec:crawl-parser}

\texttt{parser.py} thuần túy: nhận HTML (cùng URL + keyword), trả về
\texttt{<Site>JobItem}. Không có I/O, không có state, rất dễ test độc lập.
Mọi selector phải đặt trong \texttt{config.py}, không hardcode trong parser.
Parser trả \texttt{None} cho trường thiếu thay vì raise, để batch không vỡ vì
một tin lỗi. Khi site đổi markup, chỉ sửa config + parser, không đụng crawler
flow.

\texttt{utils.py} chứa các helper thuần function, không side-effect, ví dụ
\texttt{encode\_input()} chuyển \texttt{"data analyst"} thành
\texttt{"data-analyst"} cho URL slug.


\subsection{Điểm vào CLI và patch Windows}
\label{subsec:crawl-main}

\texttt{\_\_main\_\_.py} là điểm vào duy nhất khi chạy
\texttt{python -m src.crawl\_layer.crawler.<site>}.

\begin{lstlisting}[caption={Entry point CLI (src/crawl\_layer/crawler/itviec/\_\_main\_\_.py)},label={lst:main}]
async def _run(keyword, max_pages, headless) -> None:
    logging.basicConfig(level=logging.INFO, format="...")
    crawler = ItviecCrawler(keyword=keyword, max_pages=max_pages, headless=headless)
    items = await crawler.crawl()
    logging.info("Exported %d items to itviec_jobs.jsonl", len(items))

def main() -> None:
    parser = argparse.ArgumentParser(description="ITviec async crawler")
    parser.add_argument("--keyword", default="data")
    parser.add_argument("--max-pages", type=int, default=2)
    parser.add_argument("--headless", action="store_true")
    args = parser.parse_args()
    asyncio.run(_run(args.keyword, args.max_pages, headless=args.headless))

if __name__ == "__main__":
    main()
\end{lstlisting}

\textbf{Windows ProactorEventLoop patch (bắt buộc với nodriver).} Các crawler
nodriver có một khối patch ở đầu \texttt{\_\_main\_\_.py} để im lặng lỗi
\texttt{\_\_del\_\_} của \texttt{ProactorEventLoop} trên Windows (Chrome
subprocess đóng socket gây exception ồn khi thoát). Khối này không được xóa.

\begin{lstlisting}[caption={Patch Windows cho nodriver (đặt ở đầu \_\_main\_\_.py)},label={lst:winpatch}]
import sys
if sys.platform == "win32":
    from asyncio.proactor_events import _ProactorBasePipeTransport
    from asyncio.base_subprocess import BaseSubprocessTransport
    def silence_del(cls):
        orig_del = getattr(cls, "__del__", None)
        if not orig_del:
            return
        def new_del(self):
            try:
                orig_del(self)
            except Exception:
                pass
        cls.__del__ = new_del
    silence_del(_ProactorBasePipeTransport)
    silence_del(BaseSubprocessTransport)
\end{lstlisting}


\subsection{Lưu dữ liệu tạm (save\_to\_temp)}
\label{subsec:crawl-save}

Hàm duy nhất mọi crawler dùng để lưu là \texttt{save\_to\_temp()}.

\begin{lstlisting}[caption={save\_to\_temp() -- ghi append-only ra temp JSONL (src/crawl\_layer/utils/loader.py)},label={lst:save}]
def save_to_temp(data, source_name, entity_name="jobs"):
    TEMP_DIR.mkdir(parents=True, exist_ok=True)
    today_str = datetime.now().strftime("%Y%m%d")
    file_name = f"{source_name}_{entity_name}_{today_str}.jsonl"
    file_path = TEMP_DIR / file_name
    if isinstance(data, dict):
        data = [data]
    with open(str(file_path), 'a', encoding='utf-8') as f:   # mode APPEND
        for item in data:
            f.write(json.dumps(item, ensure_ascii=False) + '\n')
    return file_path
\end{lstlisting}

Đặc điểm: \textbf{append-only} (chạy lại cùng ngày cộng dồn, không ghi đè); tên
file theo ngày \texttt{<source>\_jobs\_YYYYMMDD.jsonl}; đường dẫn gốc
\texttt{TEMP\_DIR = src/crawl\_layer/temp\_data}. Việc dọn temp chỉ do Bronze
upload thực hiện.


\subsection{Kết quả dữ liệu xuất ra temp folder}
\label{subsec:crawl-temp-output}

Các dữ liệu xuất ra sẽ có dạng: \textbf{mỗi dòng trong file temp là một đối
tượng JSON thuần} (định dạng JSON Lines -- JSONL), tương ứng với đúng một tin
tuyển dụng. Vì crawler dùng \texttt{asdict()} trên dataclass
\texttt{<Site>JobItem}, mọi trường đều xuất hiện; trường nào site không cung cấp
thì mang giá trị \texttt{null}. File được lưu tại
\texttt{src/crawl\_layer/temp\_data/<source>\_jobs\_YYYYMMDD.jsonl}; ví dụ thực
tế trong dự án là \texttt{itviec\_jobs\_20260621.jsonl}.

Một dòng (đã pretty-print và rút gọn các trường text dài cho dễ đọc) có dạng
như sau:

\lstdefinestyle{jsonstyle}{
    language=,
    basicstyle=\ttfamily\footnotesize,
    showstringspaces=false,
    breaklines=true,
    frame=single,
    rulecolor=\color{lightgray},
    stringstyle=\color{teal},
    keywordstyle=\color{blue}\bfseries,
    morekeywords={null,true,false},
    columns=fullflexible,
    keepspaces=true,
}

\begin{lstlisting}[style=jsonstyle,caption={Mẫu một bản ghi JSONL xuất ra temp folder (đã rút gọn trường text dài)},label={lst:temp-sample}]
{
  "job_title": "Data Engineer - VIC",
  "company_name": "Viettel Group",
  "location": "Tang 5, Tru so chinh cua Tap doan Viettel, ..., Ha Noi",
  "job_industry": "Telecommunication",
  "job_description": "Xay dung va phat trien cac he thong du lieu lon (Big Data) ... [da rut gon]",
  "source_site": "itviec",
  "job_url": "https://itviec.com/it-jobs/data-engineer-vic-viettel-group-5243?...",
  "search_keyword": "data",
  "scraped_at": "2026-06-21T09:45:10.308085",
  "salary": "1,000 - 2,000 USD",
  "benefits": "Thu nhap canh tranh theo thi truong ... [da rut gon]",
  "requirements": null,
  "company_size": null
}
\end{lstlisting}

\textbf{Lưu ý:} trong file thật, các trường text dài (\texttt{job\_description},
\texttt{benefits}, \texttt{requirements}) giữ nguyên toàn bộ nội dung tiếng
Việt có dấu và không bị rút gọn; phần \texttt{... [da rut gon]} ở Listing~\ref{lst:temp-sample}
chỉ là cách trình bày cho vừa khung báo cáo. Toàn bộ file là \textbf{một dòng
JSON mỗi tin} (không có pretty-print), nối tiếp nhau, ghi nối tiếp (append-only)
cho đến khi Bronze upload dọn dẹp.

Các trường trong bản ghi được mô tả ở Bảng~\ref{tab:temp-fields}.

\begin{table}[h!]
\centering
\caption{Mô tả các trường trong bản ghi JSONL xuất ra temp folder}
\label{tab:temp-fields}
\begin{tabular}{p{0.22\textwidth} p{0.14\textwidth} p{0.56\textwidth}}
\toprule
\textbf{Trường} & \textbf{Kiểu} & \textbf{Mô tả} \\
\midrule
\texttt{job\_title}        & str & Tên vị trí công việc \\
\texttt{company\_name}     & str & Tên công ty tuyển dụng \\
\texttt{location}          & str & Địa điểm làm việc \\
\texttt{job\_industry}     & str & Lĩnh vực / ngành nghề \\
\texttt{job\_description}  & str & Mô tả công việc (text dài) \\
\texttt{source\_site}      & str & Tên nguồn: \texttt{itviec} / \texttt{topcv} / \texttt{vietnamworks} \\
\texttt{job\_url}          & str & URL trang chi tiết -- khóa dedup URL và upsert ở hạ lưu \\
\texttt{search\_keyword}   & str & Từ khóa tìm kiếm dùng khi cào \\
\texttt{scraped\_at}       & str & Thời điểm cào, định dạng ISO 8601 \\
\texttt{salary}            & str & Mức lương (text thô, chưa chuẩn hóa) \\
\texttt{benefits}          & str & Quyền lợi được hưởng (text dài, có thể \texttt{null}) \\
\texttt{requirements}      & str & Yêu cầu ứng viên (text dài, có thể \texttt{null}) \\
\texttt{company\_size}     & str & Quy mô công ty (trường per-site, có thể \texttt{null}) \\
\bottomrule
\end{tabular}
\end{table}

Các trường chung (\texttt{job\_title} \ldots \texttt{requirements}) đến từ
dataclass gốc \texttt{JobItem}; các trường bổ sung như
\texttt{company\_size} (ITviec) hay \texttt{job\_type},
\texttt{experience\_level}, \texttt{education\_level}, \texttt{job\_position},
\texttt{job\_deadline} (TopCV) đến từ subclass per-site. Nhờ dùng
\texttt{asdict()}, chỉ cần thêm trường vào subclass là dữ liệu tự động xuất ra
temp mà không phải sửa code ghi file.


\subsection{Tích hợp Bronze layer (pipeline hạ lưu)}
\label{subsec:crawl-bronze}

Crawl layer không tự upload S3. Sau khi crawl xong, Bronze layer đọc temp:

\begin{lstlisting}[caption={Bronze upload (src/storage\_layer/MinIO\_S3/layer/bronze/main.py)},label={lst:bronze}]
# python -m src.storage_layer.MinIO_S3.layer.bronze.main --source topcv
from src.crawl_layer.utils.loader import load_to_bronze
load_to_bronze(bronze_bucket_name, source_filter=args.source)
\end{lstlisting}

Quy trình Bronze: (1) gzip các file \texttt{<source>\_jobs\_*.jsonl}; (2) upload
\texttt{.gz} lên S3 theo key \texttt{source/<source>/...}; (3) \textbf{xóa} cả
\texttt{.gz} lẫn \texttt{.jsonl} của source đó cục bộ (destructive). Hệ quả:
\texttt{SOURCE\_NAME} phải khớp \texttt{--source} của bronze và
\texttt{bronze\_source} trong Airflow; mismatch khiến dữ liệu nằm lại temp không
được upload.


\subsection{Tích hợp Airflow DAG}
\label{subsec:crawl-airflow}

Airflow không import module crawler -- nó chỉ điều phối qua \texttt{DockerOperator}.
Mỗi site có một DAG siêu mỏng gọi factory:

\begin{lstlisting}[caption={DAG mỏng per-site (src/orchestration\_layer/dags/crawl\_topcv.py)},label={lst:dag-thin}]
from _dag_factory import create_crawl_dag
dag = create_crawl_dag("topcv")
\end{lstlisting}

Factory giữ hai bảng cấu hình phải cập nhật khi thêm source:
\texttt{SITE\_CONFIGS} (map tên site sang \texttt{crawl\_module},
\texttt{validate\_module}, \texttt{bronze\_source}, \texttt{silver\_module},
\texttt{supabase\_site}) và \texttt{CRAWL\_SCHEDULES} (cron mỗi site, offset để
không cào song song: topcv \texttt{0 */3}, itviec \texttt{15 */3}, vietnamworks
\texttt{30 */3}).

\begin{lstlisting}[caption={Cấu hình site và lệnh crawl trong DAG factory},label={lst:dag-factory}]
SITE_CONFIGS = {
    "topcv":       {"crawl_module": "topcv", "bronze_source": "topcv", ...},
    "itviec":      {"crawl_module": "itviec", "bronze_source": "itviec", ...},
    "vietnamworks":{"crawl_module": "vietnamworks", "bronze_source": "vietnamworks", ...},
}
CRAWL_SCHEDULES = {"topcv": "0 */3 * * *", "itviec": "15 */3 * * *", ...}

# Lenh crawl chay trong container (DockerOperator):
# xvfb-run ... python -m src.crawl_layer.crawler.<crawl_module>
#   --keyword {{ params.keyword }} --max-pages {{ params.max_pages }}
\end{lstlisting}

\texttt{xvfb-run} chỉ cần cho nodriver; với aiohttp nó vô hại. DAG params override
được qua Airflow UI ``Trigger DAG w/ config'': \texttt{\{"keyword": "...",
"max\_pages": N\}}.


\subsection{Quy trình thêm crawler mới}
\label{subsec:crawl-addnew}

Giả sử thêm site \texttt{careerbuzz} dùng nodriver:

\begin{enumerate}
    \item \textbf{Data model.} Thêm subclass \texttt{CareerBuzzJobItem(JobItem)}
          trong \texttt{data\_class.py}, chỉ thêm trường đặc thù, mặc định
          \texttt{None}.
    \item \textbf{Tạo package.} Tạo thư mục
          \texttt{src/crawl\_layer/crawler/careerbuzz/} với \texttt{\_\_init\_\_.py}
          (rỗng), \texttt{config.py}, \texttt{browser.py} (hoặc
          \texttt{http\_client.py}), \texttt{parser.py}, \texttt{crawler.py},
          \texttt{utils.py}, \texttt{\_\_main\_\_.py}.
    \item \textbf{config.py.} Đặt \texttt{SOURCE\_NAME = "careerbuzz"},
          \texttt{ENTITY\_NAME = DEFAULT\_ENTITY\_NAME}, URL, selector, timeout,
          env var credential, \texttt{BROWSER\_ARGS}.
    \item \textbf{parser.py.} Viết \texttt{CareerBuzzParser}, test độc lập với
          HTML sample trước.
    \item \textbf{browser.py / http\_client.py.} Sao chép từ site tương đồng
          stack, đổi selector/timing; đảm bảo \texttt{\_\_aenter\_\_}/
          \texttt{\_\_aexit\_\_} và async generator \texttt{iter\_job\_panels()}
          (nodriver) hoặc \texttt{fetch()} (aiohttp).
    \item \textbf{crawler.py.} Viết \texttt{CareerBuzzCrawler} theo pattern
          \texttt{crawl()} → per-page scrape → \texttt{\_flush\_batch()}.
    \item \textbf{\_\_main\_\_.py.} Sao chép từ \texttt{itviec/\_\_main\_\_.py}
          \textbf{bao gồm khối Windows patch} nếu dùng nodriver.
    \item \textbf{Verify cục bộ.} Chạy \texttt{python -m
          src.crawl\_layer.crawler.careerbuzz --keyword data --max-pages 1
          --headless}; kiểm tra temp file có dữ liệu.
    \item \textbf{Tích hợp Airflow.} Thêm \texttt{SITE\_CONFIGS["careerbuzz"]} và
          \texttt{CRAWL\_SCHEDULES["careerbuzz"]}; tạo file DAG
          \texttt{crawl\_careerbuzz.py} gọi \texttt{create\_crawl\_dag("careerbuzz")}.
    \item \textbf{Tích hợp Bronze.} Đảm bảo \texttt{bronze\_source = "careerbuzz"}
          khớp \texttt{SOURCE\_NAME}; thêm module validate nếu cần.
\end{enumerate}


\subsection{Các lưu ý quan trọng (gotchas)}
\label{subsec:crawl-gotchas}

\begin{enumerate}
    \item Không có \texttt{\_\_init\_\_.py} ở \texttt{src/} -- toàn bộ import
          absolute \texttt{src.*}; luôn \texttt{cd} repo root trước
          \texttt{python -m}.
    \item Temp file \textbf{append-only} -- chạy lại cùng ngày cộng dồn; dọn temp
          chỉ do Bronze upload (destructive). Dùng \texttt{--source} để cô lập
          1 site.
    \item \texttt{MinIO\_S3} là tên legacy -- thực ra nói chuyện AWS S3 qua
          \texttt{boto3} (\texttt{get\_s3\_client()}).
    \item Windows ProactorEventLoop patch ở \texttt{\_\_main\_\_.py} của mọi
          crawler nodriver -- \textbf{không xóa}.
    \item nodriver cần \texttt{xvfb-run} trên Linux/Docker (DAG factory đã wrap
          sẵn).
    \item \texttt{SOURCE\_NAME} phải nhất quán: \texttt{config.py} ↔ tên file
          temp ↔ \texttt{bronze\_source} trong Airflow ↔ \texttt{--source} của
          bronze.
    \item Selector ``xấu'' có thể cố tình -- nhiều site hash class name chống bot;
          copy y nguyên, test trang thật trước khi ``dọn dẹp''.
    \item Streaming per-page save qua \texttt{\_flush\_batch} là bắt buộc -- không
          gom hết rồi mới ghi.
    \item Dedup URL (\texttt{self.\_seen\_urls}) giữa các trang để tránh cào trùng.
    \item Airflow chỉ orchestrate -- không bao giờ import module crawler/storage;
          mọi logic chạy trong container \texttt{lakehouse-crawler} qua
          \texttt{DockerOperator}.
\end{enumerate}


\subsection{Kiểm thử}
\label{subsec:crawl-verify}

Dự án không có pytest/ruff/CI. Verify bằng cờ tiết kiệm:

\begin{lstlisting}[caption={Lệnh kiểm thử nhanh (chạy từ repo root)},label={lst:verify}]
# Crawler - gioi han 1 trang de nhanh
python -m src.crawl_layer.crawler.topcv --keyword data --max-pages 1
python -m src.crawl_layer.crawler.itviec --keyword data --max-pages 1 --headless
python -m src.crawl_layer.crawler.vietnamworks --keyword data --max-pages 1 --headless

# Kiem tra temp file
ls src/crawl_layer/temp_data/

# Silver - dry-run khong ghi
python -m src.storage_layer.MinIO_S3.layer.silver.cleaning.clean_topcv.main_process \
  --from_date 2026-06-21 --to_date 2026-06-21 --no_save

# Bronze - upload 1 source de thu
python -m src.storage_layer.MinIO_S3.layer.bronze.main --source topcv
\end{lstlisting}

Kiểm tra nhanh schema: mở file \texttt{*\_jobs\_*.jsonl}, mỗi dòng là một JSON
object; các key phải khớp trường trong \texttt{<Site>JobItem} (do dùng
\texttt{asdict()}).
